
\subsubsection{2.1 Weight Space Learning}

Neural Functional Networks (NFN; Zhou et al., 2023) introduce permutation equivariant layers for processing MLP and CNN weights. NFN leverages neuron permutation symmetries to design sample-efficient architectures. Universal Neural Functionals (UNF; Zhou et al., 2024) generalize NFN to any single architecture by automatically constructing equivariant layers for RNNs, Transformers, and arbitrary computational graphs. Deep Weight Space Networks (DWSNets; Navon et al., 2023) study equivariant architectures for deep weight spaces (e.g., Implicit Neural Representations). These methods are limited to homogeneous collections where all networks share the same architecture family.

Our hierarchical VAE extends weight space learning to heterogeneous collections via hierarchical composition. Level 1 uses architecture-specific NFN encoders (preserving local equivariance), Level 2 sacrifices neuron-level symmetries via pooling (exposing global structure), and Level 3 applies Transformer sequence modeling.

\subsubsection{2.2 Task Arithmetic and Model Editing}

Ilharco et al. (2022) demonstrate that fine-tuning directions in weight space ($\theta _finetuned - \theta _pretrained)$ encode task-specific information. These ''task vectors'' can be arithmetically combined. Ortiz-Jiménez et al. (2023) improve task arithmetic by linearizing models in tangent space. Model soup methods \citep{wortsman2022model} merge independently trained models by averaging weights. These methods require either shared base models (task arithmetic) or identical architectures (model soups). Our hierarchical VAE enables cross-architecture comparison without shared base models.

\subsubsection{2.3 Model Zoos}

ModelZooDataset (NeurIPS 2022 Dataset Track) provides large-scale collections for weight space learning research. SANE (Schürholt et al., 2024) introduces sequential weight processing for scalable model zoo learning. SANE proposes cross-architecture processing, though it is not clear from the paper whether they validate task-based clustering across architectures. Phase Transitions Model Zoo (Schürholt et al., 2025) systematically covers loss landscape phases. Our coverage audit validates that synthetic model zoos contain sufficient architecture-task diversity (77.8\% of cells with $\geq 30$ models for CNN+ResNet families). Baseline comparison to SANE is recommended future work.

